% SEGMENT OLP-0394-S01
% Part: many-valued-logic
% Chapter: three-valued-logics
% Section: lukasiewicz

\documentclass[../../../include/open-logic-section]{subfiles}

\begin{document}

\olfileid{mvl}{thr}{luk}

\olsection{லூகாசியேவிச்~(\L ukasiewicz) தருக்கம்}

பல்மதிப்புத் தருக்கத்திற்கான முதல் வெளியிடப்பட்ட, விரிவாக உருவாக்கப்பட்ட
முன்மொழிவுகளில் ஒன்றை போலந்து மெய்யியலாளர் யான் லூகாசியேவிச்~(\L
ukasiewicz) 1921 இல் வழங்கினார். அரிஸ்டாட்டிலின் கடற்போர் சிக்கல்
லூகாசியேவிச்~(\L ukasiewicz) அவர்களைத் தூண்டியது:
\emph{இன்று}, ``நாளை ஒரு கடற்போர் நிகழும்'' என்ற வாக்கியம் மெய்யும் அன்று,
பொய்யும் அன்று என்று தோன்றுகிறது; அதன் மெய்மதிப்பு இன்னும்
தீர்மானிக்கப்படவில்லை. இத்தகைய ``எதிர்கால நிச்சயமற்ற'' வாக்கியங்களுக்கு
மூன்றாவது மெய்மதிப்பை அறிமுகப்படுத்த லூகாசியேவிச்~(\L ukasiewicz)
முன்மொழிந்தார்.
\begin{quote}
அடுத்த ஆண்டின் ஒரு குறிப்பிட்ட நேரத்தில் வார்சாவில் நான் இருப்பது,
எடுத்துக்காட்டாக டிசம்பர் 21 நண்பகலில் நான் இருப்பது, தற்போதைய நேரத்தில்
நேர்முகமாகவும் எதிர்முகமாகவும் தீர்மானிக்கப்படவில்லை என்று முரண்பாடின்றி
நான் கருதலாம். ஆகவே, கொடுக்கப்பட்ட நேரத்தில் நான் வார்சாவில் இருப்பது
கூடியதே, ஆனால் கட்டாயமன்று. இந்தக் கருதுகோளின்படி, ``அடுத்த ஆண்டு டிசம்பர்
21 நண்பகலில் நான் வார்சாவில் இருப்பேன்'' என்ற கூற்று தற்போதைய நேரத்தில்
மெய்யும் அன்றாகவும், பொய்யும் அன்றாகவும் இருக்கலாம். ஏனெனில் அது இப்போது
மெய்யாக இருந்தால், எதிர்காலத்தில் நான் வார்சாவில் இருப்பது கட்டாயமாக
இருக்க வேண்டும்; இது கருதுகோளுக்கு முரணானது. மறுபுறம், அது இப்போது
பொய்யாக இருந்தால், எதிர்காலத்தில் நான் வார்சாவில் இருப்பது கூடாததாக
இருக்க வேண்டும்; இதுவும் கருதுகோளுக்கு முரணானது. ஆகவே கருதப்படும் கூற்று
இப்போது மெய்யும் அன்று, பொய்யும் அன்று; அது ``0'' அல்லது பொய்மையிலிருந்தும்
``1'' அல்லது மெய்மையிலிருந்தும் வேறுபட்ட மூன்றாவது மதிப்பைப் பெற வேண்டும்.
இம்மதிப்பை ``$\frac{1}{2}$.'' என்று நாம் குறிக்கலாம். அது ``கூடியது''
என்பதைப் பிரதிநிதித்துவப்படுத்தி, ``மெய்'', ``பொய்'' ஆகியவற்றுடன் மூன்றாவது
மதிப்பாக இணைகிறது.
\end{quote}
லூகாசியேவிச்~(\L ukasiewicz) அவர்களின் மூன்றாவது மெய்மதிப்பிற்கு
நாம்~$\Undef$ ஐப் பயன்படுத்துவோம்.\footnote{இங்கு லூகாசியேவிச்~(\L
ukasiewicz) ``கூடியது'' என்பதை இன்றைக்கு
வழக்கமில்லாத பொருளில், அதாவது கூடியதே ஆனால் கட்டாயமன்று என்ற பொருளில்,
பயன்படுத்துகிறார்.}

$\lnot$, $\land$, $\lor$ ஆகிய இணைப்பிகளுக்கான வாய்மைச் சார்புகளை இந்தப்
பொருள்கோளில் எளிதாகத் தீர்மானிக்கலாம்: எதிர்கால நிச்சயமற்ற வாக்கியத்தின்
மறுப்பும் எதிர்கால நிச்சயமற்ற வாக்கியமே; ஆகவே
$\tf{\lnot}(\Undef) = \Undef$. ஓர் இணையலின் ஓர் இணையுறுப்பு
தீர்மானிக்கப்படாமலும் மற்றது மெய்யாகவும் இருந்தால், இணையலும்
தீர்மானிக்கப்படாததே---எதிர்கால நிச்சயமற்ற இணையுறுப்பு எவ்வாறு முடிவடைகிறது
என்பதைப் பொறுத்து, இணையல் மெய்யாகவும் முடிவடையலாம், பொய்யாகவும் முடிவடையலாம்.
ஆகவே
\[
    \tf{\land}(\True, \Undef) =
\tf{\land}(\Undef, \True) =
\Undef.
\]
மற்ற இணையுறுப்பு பொய்யாக இருந்தால், இணையல் மெய்யாக முடியாது; ஆகவே
% SOURCE CORRECTION TA-MVL-007: restore the symmetric undefined/false argument order in the second conjunction equation.
\[\tf{\land}(\False, \Undef) =
\tf{\land}(\Undef, \False) = \False.\]
மற்ற மதிப்புகள்---உள்ளீடுகள் தீர்மானிக்கப்பட்ட மெய்மதிப்புகளான~$\True$
அல்லது~$\False$ ஆக இருக்கும்போது---செவ்வியல் தருக்கத்தில் உள்ளவற்றைப் போலவே
அமைகின்றன.

நிபந்தனைவாக்கியத்தைப் பொறுத்தவரை நிலைமை சற்றுச் சிக்கலானது. $q$ ஓர்
எதிர்கால நிச்சயமற்ற கூற்று என்க. $p$ பொய்யாக இருந்தால், $q$ எவ்வாறு
முடிவடைந்தாலும் $p \lif q$ மெய்யாகும்; ஆகவே
$\tf{\lif}(\False, \Undef) = \True$ என அமைக்க வேண்டும். மேலும் $p$
மெய்யாக இருந்தால், $q$ எவ்வாறு முடிவடைந்தாலும் $q \lif p$ மெய்யாகும்;
ஆகவே $\tf{\lif}(\Undef, \True) = \True$. $p$ மெய்யாக இருந்தால்,
$p \lif q$ மெய்யாகவோ பொய்யாகவோ முடிவடையலாம்; ஆகவே
$\tf{\lif}(\True, \Undef) = \Undef$. இதேபோல், $p$ பொய்யாக இருந்தால்,
$q \lif p$ மெய்யாகவோ பொய்யாகவோ முடிவடையலாம்; ஆகவே
$\tf{\lif}(\Undef, \False) = \Undef$. $p$, $q$ இரண்டும் எதிர்கால
நிச்சயமற்றவையாக இருக்கும் நேர்வு மட்டும் எஞ்சுகிறது. இந்த நோக்கத்தின்
அடிப்படையில், இந்த நேர்வில் உண்மையில்~$\Undef$ ஐயே ஒதுக்க வேண்டும்.
ஆனால் அப்படிச் செய்தால் $!A \lif !A$ ஒரு மெய்மம் \emph{அன்று}.
% SOURCE CORRECTION TA-MVL-008: repair ordinary parenthesis, agreement, and “had not trouble” prose defects.
லூகாசியேவிச்~(\L ukasiewicz) $!A \lor \lnot !A$,
$\lnot(!A \land \lnot !A)$ ஆகியவற்றைக்
கைவிடத் தயங்கவில்லை; ஆனால் $!A \lif !A$ ஐக் கைவிடத் தயங்கினார். ஆகவே அவர்
$\tf{\lif}(\Undef, \Undef) = \True$ என்று விதித்தார்.

\begin{defn}\ollabel{def:lukasiewicz}
மும்மதிப்பு லூகாசியேவிச்~(\L ukasiewicz) தருக்கம் பின்வரும் அணியைக் கொண்டு
வரையறுக்கப்படுகிறது:
\begin{enumerate}
  \item $\lnot$, $\land$, $\lor$, $\lif$ ஆகியவற்றைக் கொண்ட செந்தரக்
  கூற்று மொழி~$\Lang L_0$.
  \item மெய்மதிப்புகளின் கணம்~$V = \{\True, \Undef, \False\}$.
  \item $\True$ மட்டுமே நியமிக்கப்பட்ட மதிப்பு; அதாவது,
  $V^+ = \{\True\}$.
  \item வாய்மைச் சார்புகள் பின்வரும் அட்டவணைகளால் தரப்படுகின்றன:
  \begin{center}
    \begin{tabular}{c|c}
      $\tf{\lnot}$ & \\
      \hline
      $\True$ & $\False$ \\
      $\Undef$ & $\Undef$ \\
      $\False$ & $\True$
    \end{tabular}
    \quad
    \begin{tabular}{c|ccc}
      $\tf{\land}[\LogLuk[3]]$ & $\True$ & $\Undef$ & $\False$ \\
      \hline
      $\True$ & $\True$ & $\Undef$ & $\False$ \\
      $\Undef$ & $\Undef$ & $\Undef$ & $\False$\\
      $\False$ & $\False$ & $\False$ & $\False$
    \end{tabular}
    \\[2ex]
    \begin{tabular}{c|ccc}
      $\tf{\lor}[\LogLuk[3]]$ & $\True$ & $\Undef$ & $\False$ \\
      \hline
      $\True$ & $\True$ & $\True$ & $\True$ \\
      $\Undef$ & $\True$ & $\Undef$ & $\Undef$ \\
      $\False$ & $\True$ & $\Undef$ & $\False$
    \end{tabular}
    \quad
    \begin{tabular}{c|ccc}
      $\tf{\lif}[\LogLuk[3]]$ & $\True$ & $\Undef$ & $\False$ \\
      \hline
      $\True$ & $\True$ & $\Undef$ & $\False$ \\
      $\Undef$ & $\True$ & $\True$ & $\Undef$  \\
      $\False$ & $\True$ & $\True$ & $\True$
    \end{tabular}
  \end{center}
\end{enumerate}
\end{defn}

$\lnot$, $\land$, $\lor$ ஆகியவற்றை மட்டுமே கொண்ட எந்த
!!{formula}~$!A$ உம், அதன் எல்லா !!{propositional variable}s க்கும்
$\Undef$ ஒதுக்கப்பட்டால் மெய்மதிப்பு~$\Undef$ ஐப் பெறும் என்பதை எளிதாகக்
காணலாம். எடுத்துக்காட்டாக, $p \lor \lnot p$, $\lnot(p
\land \lnot p)$ என்ற செவ்வியல் மெய்மங்கள் $\LogLuk[3]$ இல் மெய்மங்கள்
அல்ல; ஏனெனில் $\pAssign v(p) = \Undef$ ஆகும்போதெல்லாம்
$\pValue{v}(!A) = \Undef$.

$\pAssign v(p) = \True$ அல்லது $\False$ ஆக உள்ள !!{valuation}s இல்,
$\pValue v(!A)$ அதன் செவ்வியல் மெய்மதிப்புடன் ஒத்திருக்கும்.

\begin{prop}
  $!A$ இல் உள்ள எல்லா $p$ க்கும்
  $\pAssign v(p) \in \{\True, \False\}$ எனில்,
  $\pValue v(!A)[\LogLuk[3]] = \pValue v(!A)[\LogCL]$.
\end{prop}

\begin{prob}\label{mvl:thr:luk:prob:luk-iff} $\LogLuk[3]$ இல்
  $\pValue v(!A \liff !B) = \pValue v((!A \lif !B) \land (!B \lif
  !A))$ என்று வரையறுக்கிறோம் என்க. $\liff$ இன் மெய்மை அட்டவணை என்னவாக
  இருக்கும்?
\end{prob}

பல செவ்வியல் மெய்மங்கள்~\LogLuk[3] இலும் \emph{மெய்மங்களே}; எடுத்துக்காட்டாக,
$\lnot p \lif (p \lif q)$. செவ்வியல் தருக்கத்தில் போலவே, இதை மெய்மை
அட்டவணைகளைக் கொண்டு சரிபார்க்கலாம்:
\begin{center}
\begin{tabular}{cc|cccccc}
  $p$ & $q$ & $\lnot$ & $p$ & $\lif$ & $\smash{(}p$ & $\lif$ & $q\smash{)}$ \\ \hline
  \True & \True & \False & \True & \True & \True & \True & \True \\
  \True & \Undef & \False & \True & \True & \True & \Undef & \Undef \\
  \True & \False & \False & \True & \True & \True & \False & \False \\
  \Undef & \True & \Undef & \Undef & \True & \Undef & \True & \True \\
  \Undef & \Undef & \Undef & \Undef & \True & \Undef & \True & \Undef \\
  \Undef & \False & \Undef & \Undef & \True & \Undef & \Undef & \False \\
  \False & \True & \True & \False & \True & \False & \True & \True \\
  \False & \Undef & \True & \False & \True & \False & \True & \Undef \\
  \False & \False & \True & \False & \True & \False & \True & \False \\
\end{tabular}
\end{center}

\begin{prob}
  பின்வருவன~\LogLuk[3] இல் மெய்மங்கள் என்று காட்டுக:
  \begin{enumerate}
    \item $p \lif (q \lif p)$
    \item\label{mvl:thr:luk:prob:luk-taut-2} $\lnot(p \land q) \liff (\lnot p \lor \lnot q)$
    \item\label{mvl:thr:luk:prob:luk-taut-3} $\lnot(p \lor q) \liff (\lnot p \land \lnot q)$
  \end{enumerate}
  (\olref[mvl][thr][luk]{prob:luk-taut-2},
  \olref[mvl][thr][luk]{prob:luk-taut-3} ஆகியவற்றில்,
  $!A \liff !B$ என்பதை $(!A \lif !B) \land (!B \lif !A)$ என்பதற்கான
  சுருக்கமாகக் கொள்க; அல்லது \olref[mvl][thr][luk]{prob:luk-iff} க்கான
  உங்கள் தீர்வைப் பயன்படுத்துக.)
\end{prob}

\begin{prob}
  பின்வரும் செவ்வியல் மெய்மங்கள்~\LogLuk[3] இல் மெய்மங்கள் அல்ல என்று
  காட்டுக:
  \begin{enumerate}
    % SOURCE CORRECTION TA-MVL-009: remove the unmatched closing parenthesis after q.
    \item $(\lnot p \land p) \lif q$
    \item $((p \lif q) \lif p) \lif p$
    \item $(p \lif (p \lif q)) \lif (p \lif q)$
  \end{enumerate}
\end{prob}

எல்லாச் செவ்வியல் மெய்மங்களும்~$\LogLuk[3]$ இல் மெய்மங்கள் அல்ல என்றாலும்,
அவை ஒவ்வொரு !!{valuation} இலும் குறைந்தபட்சம்~$\True$ அல்லது~$\Undef$
மதிப்பைப் பெற வேண்டும் என்று ஒருவர் எண்ணலாம். இதுவும் சரியன்று. ஒரு
எதிரெடுத்துக்காட்டு:
\[
  \lnot(p \lif \lnot p) \lor \lnot(\lnot p \lif p)
\]
$p$ ஆனது~$\Undef$ ஆக இருந்தால் இது~$\False$.

\begin{prob}
  பின்வரும் தொடர்புகளில் எவை லூகாசியேவிச்~(\L ukasiewicz) தருக்கத்தில்
  பொருந்துகின்றன?
  ஒவ்வொன்றுக்கும் ஒரு மெய்மை அட்டவணை தருக.
  \begin{enumerate}
    \item $p, p \lif q \Entails q$
    \item $\lnot\lnot p \Entails p$
    \item $p \land q \Entails p$
    \item $p \Entails p \land p$
    \item $p \Entails p \lor q$
  \end{enumerate}
\end{prob}

தமது மும்மதிப்பு முறைமையின் அடிப்படையில் கூடுதன்மைக்கான தருக்கம் ஒன்றைக்
கட்டமைக்க லூகாசியேவிச்~(\L ukasiewicz) விரும்பினார். அதற்காக
``$!A$ கூடியது'' என்பதற்கான
ஓரிட இணைப்பி~$\Diamond !A$ ஐயும், அதற்குரிய ``$!A$ கட்டாயமானது''
என்பதற்கான~$\Box !A$ ஐயும் அறிமுகப்படுத்தினார்:
\begin{center}
  \begin{tabular}{c|c}
    $\tf{\Diamond}$ & \\
    \hline
    $\True$ & $\True$ \\
    $\Undef$ & $\True$ \\
    $\False$ & $\False$
  \end{tabular}
  \quad
  \begin{tabular}{c|c}
    $\tf{\Box}$ & \\
    \hline
    $\True$ & $\True$ \\
    $\Undef$ & $\False$ \\
    $\False$ & $\False$
  \end{tabular}
\end{center}
வேறு சொற்களில், $p$ ஏற்கெனவே பொய் என்று தீர்மானிக்கப்படவில்லை எனில்,
அப்போதும் அப்போது மட்டுமே, அது கூடியது; $p$ ஏற்கெனவே மெய் என்று
தீர்மானிக்கப்பட்டிருந்தால், அப்போதும் அப்போது மட்டுமே, அது கட்டாயமானது.

\begin{prob}
  $\Box$, $\Diamond$ ஆகியவற்றின் மெய்மை அட்டவணைகளால் விரிவாக்கப்பட்ட
  ~\LogLuk[3] இல் $\Box p \liff \lnot\Diamond \lnot p$, $\Diamond p \liff
  \lnot \Box \lnot p$ ஆகியவை மெய்மங்கள் என்று காட்டுக.
\end{prob}

% SOURCE CORRECTION TA-MVL-010: remove repeated “However” and restore the spelling of possibility.
ஆனால் இம்முன்மொழியப்பட்ட மாதிரி தருக்கத்தின் குறைபாடுகள் விரைவில்
வெளிப்பட்டன: நிகழ்வுகள் எவ்வாறு முடிந்தாலும், $p \land \lnot p$ ஒருபோதும்
மெய்யாக முடியாது. எனவே அது இப்போது தீர்மானிக்கப்படாமல் (அதனால்
தீர்மானமற்றதாக) இருந்தாலும், கூடாததாகக் கணக்கிடப்பட வேண்டும்; அதாவது,
$\lnot \Diamond(p \land \lnot p)$ ஒரு மெய்மமாக இருக்க வேண்டும். ஆனால்
$\pAssign v(p) = \Undef$ எனில்,
$\pValue v(\lnot \Diamond(p \land \lnot p)) = \Undef$. கூடுதன்மை,
கட்டாயத்தன்மை போன்ற மாதிரி வேறுபாடுகளுக்கு இரண்டு மெய்மதிப்புகள் போதாது
என்பதில் லூகாசியேவிச்~(\L ukasiewicz) சரியாக இருந்தாலும், மூன்றாவது மெய்மதிப்பை
அறிமுகப்படுத்துவதும் போதாது.

\end{document}
